%%%%%%%%%%%%%%%%%%%%%%%%%%%%%%%%%%%%%%%%%%%%%%%%%%%%%%%%%%%%%%%%%%%%%%%%%%%%%%%%%%%%%%%%%%
% Theoretical motivations
%%%%%%%%%%%%%%%%%%%%%%%%%%%%%%%%%%%%%%%%%%%%%%%%%%%%%%%%%%%%%%%%%%%%%%%%%%%%%%%%%%%%%%%%%%
NEW TEXT MAYBE?


\subsection{Baryon asymmetry}

While energy can be converted freely into matter and anti-matter our stable world is made of
matter, where we ignore anti-matter produced in $\beta$-decay, cosmic ray showers or experimental
apparati. The evidence also favors the hypothesis that the bulk of the universe consists of
the same charge-conjugation of matter as our solar system~\cite{Coppi:2004za,Steigman:1976ev}
If the visible matter in the universe is of the same charge-handedness as our immediate world,
we are left with two possibilities. First, that the universe started out in an asymmetric state.
However, it is predicted that the constituents would quickly evolve into an symmetric state
in the hot, early universe.
The second hypothesis, that the universe started out in a symmetric state but evolved into an
asymmetric state. This raises the question of how the universe evolved from an early neutral
universe consisting of gauge bosons, fermions and quarks to the current neutral consisting
universe of predominantly hydrogen? To more quantitatively discuss the current asymmetry,
a metric is needed to measure how large this asymmetry is. Let $B$ represent 
a baryon right now and not a $B$ meson and $N_B$ the number of baryons. 
We can then talk about the $B + \bar{B}$ annihilations in
the early universe which produced two photons for every annihilation. This is used to define
an asymmetry parameter, $\eta$.

\begin{eqnarray*}
    \eta &=& \frac{N_B + N_{\bar{B}}}{N_\gamma} \\
         &\approx& 10^{-9}
\end{eqnarray*}
 


\subsection{Sakharov conditions}

In 1967, Andrei Sakharov put forth the three conditions required to achieve the matter
(baryon) asymmetry observed today, assuming a symmetric initial state in the paper “Vi-
olation of CP Invariance, c Asymmetry, and Baryon Asymmetry of the Universe.” These
conditions have become known as the Sakharov conditions~\cite{Sakharov:1967dj}
One of the necessary conditions was that there must exist a process which violates baryon number. Baryon number
violation actually does exist in the Standard Model. in the form of sphalerons, a non-
perturbative process. This is predicted to occur at very high temperatures, T = 100GeV
($10^{15}$ K), far beyond the reach of any conceivable experiment.
One of the other Sakharov conditions is CP -violating processes. These processes have been
measured in kaon and B decays, but the amount of CP -violation observed when combined
with the predicted contributions from sphalerons is 10 orders of magnitude too small! The
conclusion is that there are still unobserved processes which violate CP to a greater degree
or there is an unobserved baryon number violating process.


Sakharov mentions another possible processes in his 1967 paper, a vector boson with charge=$\pm\frac{1}{3}$,
$\pm\frac{2}{3}$ or $\pm\frac{4}{3}$,
that mixes quark and lepton flavor with M=10-1000 GeV/c$^2$ . In 1980 in his Nobel
prize acceptance speech, James Cronin used this picture to illustrate proton decay, 
referring to the mediating boson as an X-boson. This coupling also appears in SU(5) and other
GUT’s. These theories also give rise to Y-bosons which are equated with scalar Higgs
states. These decays are referred to as $B-L$ processes, where the difference of baryon
number and lepton number are conserved, though each is individually violated.
In the first SU(5) paper~\cite{Georgi:1974sy}, it was shown that proton decay was an implication of this theory.
But the experimental bounds on the lifetime of the proton is greater than 10$^{32}$ years! The
mediating boson for this process has fractional charge: X is $q=\frac{1}{3}$ . 


Other processes could violate baryon number by 2 units, for example neutron oscillations in 
which a neutron oscillates into an anti-neutron. Searches for this interaction
have produced null results~\cite{Baldo-Ceolin:1994hzw, Super-Kamiokande:2020bov}.\\

We propose a search for mesons that decay to two baryons, in our specific case

$$B^+ \rightarrow p \Lambda^0$$

though we point out that we do not tag the other-side $B$ and so the initial state
could also be a $\overline{B^0}$. Any result could be incorporated into an effective
field theory, however for this analysis we are focused on the experimental result
and leave theoretical interpretations to others.
